\section{New prospectuses: payment, scope and external law}
\label{sec:prospectus-difficulty}

A bond can promise repayment at its principal amount while another provision
allows an authority to reduce that claim. Equally, a document can describe a
reduction of outstanding principal because the investor has already received
part of the money. These sentences can look similar to a text classifier, but
they describe different changes in the investor's rights. The distinction
must be settled before the resulting product classification is used in a
selling decision.

On 4 October 2026, the unchanged prospectus reader was exercised on 24 newly
acquired public PDFs. Twenty were issue-specific documents and four were
programme or base prospectuses. They represented 16 issuer families; several
documents covered multiple series, and the Ecuador and Volcan pairs included
an initial offering and a later tap of the same series. The reader abstained
on 21 documents and returned three conditional positive readings. An abstention
means that its premises remained unresolved, not that the bond was found to
lack a loss mechanism. One positive used demonstrably incorrect supporting
clauses. Appendix~\ref{app:prospectus-difficulty} identifies every document,
its source date, retained PDF and observed result.

\subsection{A smaller outstanding balance need not be a loss}

Mizuho Markets Cayman's August 2026 programme provides the clearest error.
Its form of final terms says that amortizing notes are partially redeemed at
the specified Amortization Amount and that their outstanding principal is
reduced accordingly. A neighbouring zero-coupon option likewise reduces
outstanding principal \emph{upon payment} of the amortization amount. Both
passages became principal-write-down witnesses in the reader's positive
answer.\footnote{\href{../prospectus/difficulty-2026-10-04/pdfs/case-105635287.pdf}
{Mizuho Markets Cayman LP, US\$10 billion medium-term note programme,
31 August 2026}, PDF pp.~125--126 and 128.}

The payment relation explains why that inference is wrong. In a simple
illustration, an investor owed 100 receives 20 and retains a claim for 80;
the reduction records repayment. If 20 of an unpaid claim is extinguished,
the investor instead retains 80 without having received the missing 20.
The programme passage describes the first relation. Recognising the verb
\emph{reduced} without the payment that causes the reduction loses the
distinction. Moreover, these passages are unselected form options.
Correctly interpreting one option would still not establish which terms
govern a particular issued bond.

Ethniki's September 2026 restricted Tier 1 notes supply a contrasting
principal-loss clause. Their Prevailing Principal Amount may be Written Down,
and redemption can occur at that reduced amount without a later entitlement
to the difference. The reader found this mechanism but left 57 candidate
records unresolved. MUFG's August 2026 base prospectus similarly states that
the principal of \emph{Subordinated Notes} is written down to zero following
a Non-Viability Event. That clause supports a reading of the subordinated
option; it does not classify every note issued under the programme.
\footnote{\href{../prospectus/difficulty-2026-10-04/pdfs/case-105585556.pdf}
{Ethniki, EUR 100 million 6.875 percent RT1 notes, 7 September 2026},
PDF pp.~44--45;
\href{../prospectus/difficulty-2026-10-04/pdfs/case-105541740.pdf}
{MUFG, US\$50 billion programme, 7 August 2026}, PDF pp.~1, 20 and 88.}

\subsection{The principal must belong to the right instrument}

A provision's subject matters as much as its action. All three unresolved
records in EDF's 2023 perpetual-note document discussed separate OCEANE
convertible bonds due in 2024. An apparent exclusion of principal write-down
elsewhere in the same document applied to \emph{Exchanged Notes or Varied
Notes}. It was therefore unsuitable as an exhaustive negative statement
about the original notes. Volcan's 2025 memorandum supplied an even sharper
contrast: several candidate forfeiture clauses concerned mining concessions
or water licences. Those assets can affect the issuer's financial position,
but forfeiture of a mining concession is not itself cancellation of the
principal of its Notes.\footnote{
\href{../prospectus/difficulty-2026-10-04/pdfs/case-103648885.pdf}
{EDF, US\$1.5 billion perpetual notes, 15 June 2023}, PDF pp.~15, 33, 49,
57 and 136;
\href{../prospectus/difficulty-2026-10-04/pdfs/case-105062309.pdf}
{Volcan, US\$750 million senior secured notes, 14 November 2025},
PDF pp.~35, 80 and 116--117.}

Missing context can also produce a misleading impression of completeness.
Zurich Finance Ireland's eight-page July 2026 pricing supplement identifies
senior notes and a redemption amount, but expressly incorporates the
19 May 2026 base prospectus. The isolated-document reader left debt status
unknown while marking its local clause coverage complete. Both PDFs had
been downloaded, but the diagnostic input had not declared their dependency.
The useful repair is to connect the correct issue documents and interpret
their labelled fields. The absence of a clause from a short supplement
cannot establish that the complete terms contain no such clause.
\footnote{\href{../prospectus/difficulty-2026-10-04/pdfs/case-105489584.pdf}
{Zurich Finance (Ireland) DAC, Series 70, tranche 1, pricing supplement
of 8 July 2026}, PDF pp.~1--3;
\href{../prospectus/difficulty-2026-10-04/pdfs/case-105392707.pdf}
{Zurich EMTN base prospectus, 19 May 2026}.}

\subsection{The applicable legal power needs a date and a subject}

HPB's July 2026 senior preferred notes illustrate the next step beyond a
repayment promise. The disclosure says that an authority may order a
write-down of the Notes or convert them into CET1 instruments, and describes
cancellation of liabilities including the Notes. The reader left these
clauses unresolved. The source supports investigation of a principal-loss
power; the phrase \emph{CET1 instruments} still needs interpretation before it
can establish the narrower common-share conversion feature. An exclusion
of contractual loss terms in the definition of replacement Qualifying
Securities does not negate statutory powers over the original
Notes.\footnote{\href{../prospectus/difficulty-2026-10-04/pdfs/case-105490784.pdf}
{Hrvatska Poštanska Banka, EUR 150 million senior preferred notes,
9 July 2026}, PDF pp.~28, 31 and 59.}

Insurance-resolution language shows why the date cannot be inferred from
the presence of a legal reference. AXA's June 2024 prospectus calls the
Insurance Recovery and Resolution Directive proposal a proposed measure
and says \emph{If adopted}. SCOR's June 2026 prospectus describes national
implementation and application in January 2027 and states that French
implementation was pending. These passages disclose a possible future
change. They do not, by themselves, establish an already applicable power
at the relevant assessment date. Existing contractual provisions and other
national powers must be examined separately.\footnote{
\href{../prospectus/difficulty-2026-10-04/pdfs/case-104224021.pdf}
{AXA, EUR 125 million senior notes, 12 June 2024}, PDF pp.~11--12;
\href{../prospectus/difficulty-2026-10-04/pdfs/case-105421783.pdf}
{SCOR, EUR 500 million subordinated notes, 2 June 2026}, PDF pp.~27--28.}

Nor is the question confined to regulatory capital. Prudentia's AirBlu
compartment terms extinguish the right to further amounts after the
compartment assets have been realised and distributed. Orange describes
court restructuring capable of binding dissenting creditor classes. Whether
these mechanisms satisfy the declared principal-loss predicate requires
analysis of the unpaid claim and the authority for the change. Neither
can be excluded merely by calling it ordinary credit risk or observing
that the security is not bank capital.\footnote{
\href{../prospectus/difficulty-2026-10-04/pdfs/case-103216017.pdf}
{Prudentia AirBlu, EUR 35 million notes, 13 October 2022}, PDF p.~44;
\href{../prospectus/difficulty-2026-10-04/pdfs/case-105155475.pdf}
{Orange, US\$6 billion notes, 13 January 2026}, PDF p.~45.}

The 24-document exercise therefore identifies several different engineering
tasks: interpret payment and extinguishment correctly; bind a clause to its
instrument; assemble the issue documents; and establish which law applies
when. The sources and results remain preserved with the original reader
unchanged. The findings are based on inspected extracted text and are
provisional pending rendered-source and independent legal review. Passing
software tests cannot supply those missing interpretations.


\subsection{Three further historical documents}

A further search retained three historical documents for testing the boundary
between contractual promises and local-law investigation. Banco Espírito
Santo's final terms dated 6 May 2014 identify EUR 750 million of 2.625 percent
senior notes due 8 May 2017, Series 36, tranche 1. Page 2 promises redemption
at 100 percent of nominal amount, subject to the stated prior-redemption and
cancellation qualifications. Page 1 incorporates the July 2013 prospectus and
four dated supplements. Thus the visible senior-status and repayment fields
are precise contractual evidence, but the five-page final terms alone cannot
answer a question about later Portuguese resolution measures or their
recognition abroad.\footnote{\href{../prospectus/jurisdiction-2026-10-04/sources/jur-bes-2014-ptbeqkom0019/source.pdf}
{BES, Series 36, tranche 1, final terms of 6 May 2014}, PDF pp.~1--4.
Page 4 confirms ISIN PTBEQKOM0019. The retained PDF is a scan; these
fields were inspected on rendered pages.}
The proposed follow-up is to match the exact security to the relevant
Portuguese measures and judicial decisions. That matching has not yet been
established by the retained external-law sources.

The Greek document gives a directly visible example of a broad promise
qualified by another jurisdiction's law. Its March 2008 offering circular
covers a EUR 600 million tap of inflation-linked notes due in 2057.
Condition 4 floors the contractual redemption amount at principal.
Condition 14 selects English law and contains a broad waiver of immunity,
but then expressly subjects execution and attachment of the Republic's
property to international conventions and Greek law. A repayment entitlement,
jurisdiction over a lawsuit and access to assets for enforcement are therefore
three different propositions. The country name also cannot justify attaching
a Greek-law restructuring rule to these English-law notes without checking
that rule's scope and the actual series history.\footnote{
\href{../prospectus/jurisdiction-2026-10-04/sources/jur-hellenic-2008-candidate/source.pdf}
{Hellenic Republic, 2.085 percent inflation-linked notes due 25 July 2057,
offering circular of 28 March 2008}, PDF pp.~10 and 17--18.
The permanent ISIN is XS0292467775 (PDF p.~25).}
No specific seizure, restructuring outcome or recovery amount is established
by that qualification alone.

The Slovenian search produced SID Bank's March 2011 memorandum, rather than
an NLB subordinated issue. It describes a EUR 350 million tap of 3 percent
notes due in 2015 and an unconditional, irrevocable state guarantee. Its
description of the guarantee refers to article 13 of the Slovenian Export
and Development Bank Act, transactions under articles 11--12 and a creditor's
written request. The Notes themselves use English law. The correct
investigation must therefore establish the statutory guarantee's scope and
conditions, while keeping SID Bank distinct from another Slovenian bank
with different instruments or intervention history.\footnote{
\href{../prospectus/jurisdiction-2026-10-04/sources/jur-nlb-2011-candidate/source.pdf}
{SID Bank, EUR 350 million 3 percent notes due 2015, information memorandum
of 9 March 2011}, PDF pp.~1, 18 and 24. The retained filename records the
original search candidate; the document's issuer is SID Bank.}
This document is retained as an issuer-and-guarantee control. It does not
demonstrate that the guarantee was overridden or that these notes suffered
a principal write-down.

These new documents extend the tests in two directions. The BES terms
provide a concrete repayment promise requiring a later-event investigation.
The Greek and SID documents prevent that investigation from becoming a
country-wide presumption: governing law, enforcement limits, issuer identity
and the statutory guarantee each have to be bound to the particular claim.
The court-source acquisitions attempted for Portuguese, Greek and Slovenian
litigation had not yielded readable judgments at this checkpoint, so the
corresponding holdings and event outcomes remain open.
